\TieudeT{PHONG TỎA VẬT LÍ}
\subsubsection*{PHẦN I. Thí sinh trả lời câu hỏi từ câu 1 đến câu 18. Mỗi câu hỏi thí sinh chỉ chọn một phương án}
\setcounter{ex}{0}
\begin{ex}%book
Hệ số ma sát trượt giữa vật với mặt tiếp xúc là $\mu$, phản lực mà mặt tiếp xúc tác dụng lên vật là $N$. Lực ma sát trượt tác dụng lên vật là $F_{mst}$. Hệ thức đúng là
\choice
{$F_{mst}=\dfrac{N}{\mu_{t}}$}
{$F_{mst}=\mu_{t}.N^2$}
{$F_{mst}=\mu_{t}^2.N$}
{\True $F_{m s t}=\mu_{t}.N$}
\loigiai{

Biểu thức tính độ lớn lực ma sát trượt: $F_{m s t}=\mu_{t} N$

Trong đó:

$+\mu_{t}$ là hệ số ma sát trượt giữa vật và mặt tiếp xúc

\begin{itemize}
\item $N$ là áp lực của vật lên mặt tiếp xúc
\end{itemize}
}
\haidong{10} \end{ex}
\begin{ex}%book
Phát biểu nào sau đây là \textbf{sai}? Khi một vật trượt trên một mặt phẳng, độ lớn của lực ma sát trượt
\choice
{\True phụ thuộc vào diện tích tiếp xúc của vật với mặt phẳng đó}
{không phụ thuộc vào tốc độ của vật}
{tỉ lệ với độ lớn áp lực của vật lên mặt phẳng đó}
{phụ thuộc vào vật liệu và tình trạng của hai mặt tiếp xúc}
\loigiai{
Lực ma sát trượt
\begin{itemize}
\item phụ thuộc vào tình trạng, chất liệu của hai mặt tiếp xúc;
\item tỉ lệ với áp lực của vật lên mặt tiếp xúc;
\item không phụ thuộc vào tốc độ của vật;
\item không phụ thuộc vào diện tích tiếp xúc giữa hai mặt tiếp xúc.
\end{itemize}
}
\haidong{10} 
\end{ex}
\begin{ex}%book%[1P2B7-1][Loai1]
Một vật đặt trên mặt bàn nằm ngang, có diện tích tiếp xúc là $S$. Trong một phạm vi khá rộng, khi diện tích tiếp xúc tăng gấp đôi thì lực ma sát trượt xuất hiện giữa vật và mặt tiếp xúc khi vật chuyển động sẽ
\choice
{tăng gấp đôi}
{\True gần như được giữ không đổi}
{giảm một nửa}
{phụ thuộc diện tích mặt tiếp xúc $S$}
\loigiai{
Độ lớn lực ma sát trượt không phụ thuộc vào diện tích tiếp xúc.
}
\haidong{10} \end{ex}
\begin{ex}%book%[1P2Y7-1][Loai1]
Một vật trượt trên một mặt phẳng, khi tốc độ của vật tăng thì hệ số ma sát giữa vật và mặt phẳng
\choice
{\True không đổi}
{tăng tỉ lệ với bình phương tốc độ của vật}
{giảm xuống}
{tăng tỉ lệ với tốc độ của vật}
\loigiai{
Hệ số ma sát trượt giữa vật và mặt phẳng không phụ thuộc vào tốc độ trượt của vật trên mặt phẳng nghiêng đó.
}
\haidong{10} \end{ex}
\begin{ex}%book
Vai trò của lực ma sát trượt \textbf{không} được thể hiện trong hiện tượng nào sau đây?
\choice
{Phanh xe đạp, ô tô, xe máy giúp xe dừng lại}
{Người ta quẹt que diêm vào vỏ hộp diêm để tạo ra lửa}
{Vào mùa đông, xoa hai bàn tay vào nhau giúp tay ta ấm lên}
{\True Quyển sách đang nằm yên trên mặt phẳng nghiêng}
\loigiai{
Lực ma sát nghỉ giữ cho quyển sách nằm yên trên mặt phẳng nghiêng.
}
\haidong{10} \end{ex}
\begin{ex}%book
Một vật trượt có ma sát trên một mặt tiếp xúc nằm ngang. Nếu diện tích tiếp xúc của vật đó giảm $3$ lần thì độ lớn lực ma sát trượt giữa vật và mặt tiếp xúc sẽ
\choice
{giảm $3$ lần}
{tăng $3$ lần}
{giảm $6$ lần}
{\True không thay đổi}
\loigiai{
Lực ma sát trượt không phụ thuộc vào diện tích tiếp xúc giữa hai mặt tiếp xúc nên khi diện tích tiếp xúc tăng lên hay giảm đi thì độ lớn lực ma sát trượt không thay đổi.
}
\haidong{10} \end{ex}
\begin{ex}%book%[1P2B7-1][Loai1]
Phát biểu nào sau đây là đúng khi nói về đặc điểm của lực ma sát nghỉ?
\choice
{\True Xuất hiện khi một vật chịu tác dụng của ngoại lực có xu hướng làm cho vật chuyển động nhưng vẫn đứng yên}
{Lực ma sát nghỉ luôn nhỏ hơn ngoại lực tác dụng vào vật}
{Lực ma sát nghỉ tỉ lệ với áp lực N của vật lên mặt phẳng đỡ}
{Lực ma sát nghỉ luôn vuông góc với bề mặt tiếp xúc}
\loigiai{
Lực ma sát nghỉ xuất hiện khi vật chưa chuyển động.
}
\haidong{10} \end{ex}
\begin{ex}%book%[1P2B7-1][Loai1]
Một vật đặt trên mặt bàn nằm ngang, tác dụng vào vật một lực $F$ có phương song song với mặt tiếp xúc, càng tăng lực tác dụng đặt vào vật sao cho vật vẫn đứng yên, trong trường hợp này thì lực ma sát nghỉ có giá trị
\choice
{không đổi}
{\True càng tăng theo}
{càng giảm dần}
{tăng lên rồi giảm}
\loigiai{
Lực ma sát nghỉ có giá trị tăng theo ngoại lực khi vật vẫn đứng yên.
}
\haidong{10} \end{ex}
\begin{ex}%book
Phát biểu nào sau đây là \textbf{sai}? Khi một người đi bộ, lực ma sát tác dụng lên chân người đó
\choice
{là lực ma sát nghỉ}
{\True có hướng ngược với hướng chuyển động của người}
{có phương song song với mặt đường}
{có vai trò giúp người đó tiến về phía trước}
\loigiai{
Khi người đi bộ, chân người đó tác dụng lên mặt tiếp xúc một lực hướng về phía sau. Mặt tiếp xúc sẽ tác dụng lên chân người đó một lực ma sát nghỉ có hướng ngược lại, như vậy, lực này có hướng cùng hướng chuyển động của người và có vai trò là lực phát động giúp người tiến về phía trước.
}
\haidong{10} \end{ex}
\begin{ex}%book %[1P2B7-1][Loai1]
Phát biểu nào sau đây là \textbf{sai} khi nói về lực ma sát?
\choice
{Lực ma sát trượt chỉ xuất hiện khi vật này trượt trên bề mặt của vật khác}
{Hướng của lực ma sát trượt ngược với hướng chuyển động của vật}
{Lực ma sát lăn tỉ lệ với áp lực do vật tác dụng lên mặt tiếp xúc}
{\True Viên gạch nằm yên trên mặt phẳng nghiêng không chịu tác dụng của lực ma sát}
\loigiai{ 
Viên gạch nằm yên trên mặt phẳng nghiêng không chịu tác dụng của lực ma sát là \textbf{sai} vì nó chịu tác dụng của lực ma sát nghỉ.
}
\haidong{10} \end{ex}
\begin{ex}%book%[1P2B7-1][Loai1]
Một trái bida đỏ lúc đầu đứng yên trên bàn bida (mặt phẳng nằm ngang, nhám). Sau khi cơ thủ đánh, trái bida đỏ được truyền một vận tốc đầu, nó chuyển động chậm dần do
\choice
{\True lực ma sát}
{phản lực}
{lực quán tính}
{lực đàn hồi}
\loigiai{
Viên bi chuyển động chậm dần do lực ma sát.
}
\haidong{10} \end{ex}
\begin{ex}%Tailieuchuan
Một xe có khối lượng $m=5$ tấn đang đứng yên trên mặt phẳng nghiêng $30^\circ$ so với mặt ngang. Độ lớn của lực ma sát tác dụng lên xe
\choice
{lớn hơn trọng lượng của xe}
{bằng trọng lượng của xe}
{bằng thành phần trọng lực vuông góc với mặt phẳng nghiêng}
{\True bằng thành phần trọng lực song song với mặt phẳng nghiêng}
\loigiai{}
\end{ex}
\begin{ex}%Tailieuchuan
\immini[thm]{Một thùng gỗ được kéo bởi lực $\vec{F}$ như hình vẽ. Thùng chuyển động thẳng đều. Công thức xác định lực ma sát nào sau đây là đúng?
\choice
{\True $F_{ms}=F\cos \alpha$}
{$F_{ms}=P$}
{$F_{ms}=F \sin \alpha$}
{$F_{ms}=F \cos \alpha-P$}}{\begin{tikzpicture}[scale=1,thick,>=stealth',draw=Mapcolor]
\fill[pattern=north east lines] (0,-0.2) rectangle (6,0);
\draw (0,0)--(6,0);
\fill[pattern=north west lines,draw] (2,0) rectangle (3.5,0.8);
\draw[->,very thick,Mapcolor](3.5,0.4)--++(30:2)node[right]{$\vec F$};
\draw[densely dashed] (3.5,0.4)--(6,0.4);
\draw (3.5,0.4)++(0:0.6) arc (0:30:0.6);
\path(3.5,0.4)--++(15:0.8) node{{$\alpha$}};
\end{tikzpicture}}
\loigiai{}
\end{ex} 
\begin{ex}%book%[1P2K7-2][Loai1]
Biết hệ số ma sát trượt của vật và sàn là $0,2$. Cần kéo một vật trọng lượng $20\ N$ với một lực theo phương ngang có độ lớn bằng bao nhiêu để vật chuyển động đều trên một mặt sàn nằm ngang?  
\choice
{\True $4\ N$}
{$10\ N$}
{$2\ N$}
{$5\ N$}
\haidong{3}\loigiai{
Để vật chuyển động thẳng đều thì $F = F_{mst} = \mu mg = 0,2.20 = 4\ N$.
}
\haidong{10} \end{ex}
\begin{ex}%book%SBT KNTT 18.6
Một toa tàu có khối lượng $60$ tấn chuyển động thẳng đều dưới tác dụng của lực kéo của đầu tàu theo phương nằm ngang $F=4,5.10^4\ N$. Lấy $g=10\ m/s^2$. Hệ số ma sát giữa tàu và đường ray là
\choice
{\True $0,075$}
{$0,06$}
{$0,15$}
{$0,015$}
\haidong{4}\loigiai{
\begin{itemize}
\item Vì toa tàu chuyển động thẳng đều nên: $F_{\ct{kéo}}=F_{ms}\Rightarrow \mu mg=4,5.10^4\Rightarrow \mu=0,075$.
\end{itemize}
}
\haidong{10} \end{ex}
\begin{ex}%book %[1P2K7-2][Loai1]
Một khối gỗ có khối lượng $4\ kg$ nằm trên mặt bàn. Hệ số ma sát nghỉ giữa vật và mặt bàn là $0,25$. Lấy $ g =10\ m/s^2$. Muốn vật chuyển động được trên mặt bàn thì lực kéo tác dụng vào vật phải có độ lớn tối thiểu là 
\choice
{$4\ N$}
{\True $10\ N$}
{$80\ N$}
{$20\ N$}
\haidong{4}\loigiai{ 
Để vật chuyển động thì lực kéo tác dụng lên vật phải lớn hơn hoặc bằng lực ma sát nghỉ $F_k \geq F_{msn} = \mu_n.mg = 0,25.4.10 = 10$ N.
}
\haidong{10} 
\end{ex}
\begin{ex}%book
Cho một vật có khối lượng $10\ kg$ đặt trên một sàn nhà. Một người tác dụng một lực là $30\ N$ kéo vật theo phương ngang, hệ số ma sát giữa vật và sàn nhà là $\mu=0,2$. Cho $g=10\ m/s^2$. Gia tốc của vật có độ lớn là
\choice
{$4\ m/s^2$}
{$3\ m/s^2$}
{$2\ m/s^2$}
{\True $1\ m/s^2$}
\loigiai{

Chọn hệ quy chiếu như hình vẽ vật chịu tác dụng của các lực: $\overrightarrow{N}, \overrightarrow{P}, \overrightarrow{F_{ms}}, \overrightarrow{F}$

Theo định  II Newton ta có: $\overrightarrow{N}+\overrightarrow{P}+\overrightarrow{F_{ms}}+\overrightarrow{F}=ma$

Chiếu lên trục $Ox: F-F_{ms}=$ ma (1)

Chiếu lên trục $Oy: N-P=0 \Rightarrow N=P=mg$ (2)

Từ (1) và (2):

$
a=\dfrac{F-\mu m g}{m}=\dfrac{30-0,2 \cdot 10.10}{10}=1\ m/s^2
$

\begin{center}
\includegraphics[width=4cm]{images/2023_11_19_be4543c46a9f67e78874g-12(1)}
\end{center}

Chọn D
}
\haidong{10} \end{ex}
\begin{ex}%book%SBT KNTT 18.8
Một ô tô có khối lượng $1,5$ tấn chuyển động trên đường nằm ngang. Hệ số ma sát của xe với mặt đường là $0,01$. Biết lực kéo gây ra bởi động cơ song song với mặt đường. Lấy $g=10\ m/s^2$. Độ lớn của lực kéo để ô tô chuyển động nhanh dần đều với gia tốc $0,25\ m/s^2$ là	
\choice
{\True $525\ N$}
{$255\ N$}
{$552\ N$}
{$545\ N$}
\haidong{6}\loigiai{
\begin{itemize}
\item $F=ma+\mu mg=1500.(0,25+0,01.10)=525\ N$.
\end{itemize}
}
\haidong{10} 
\end{ex}
\subsubsection*{PHẦN 2. Thí sinh trả lời từ câu 1 đến câu 4. Trong mỗi ý a), b), c), d) ở mỗi câu, thí sinh chọn đúng hoặc sai.}
\setcounter{ex}{0}
\begin{ex}%book
Dưới tác dụng của một lực kéo không đổi song song với mặt tiếp xúc, một viên gạch hình hộp chữ nhật trượt trên một tấm ván khô được đặt nằm ngang. Để lực ma sát trượt thay đổi thì ta thực hiện cách làm nào sau đây?
\choiceTF[t]
{\True Làm ướt tấm ván}
{\True Nâng tấm ván lên thành một mặt phẳng nghiêng}	
{Thay đổi tốc độ chuyển động của vật}
{\True Lật viên gạch sang một mặt tiếp xúc bằng chất liệu khác}	
\loigiai{
\begin{itemchoice}
\itemch
\itemch
\itemch
\itemch	
\end{itemchoice}
}
\end{ex}

\begin{ex}%book
Trong đời sống và kĩ thuật, lực ma sát có thể có lợi hoặc có hại tùy trường hợp cụ thể. Các nhận định sau liên quan đến ứng dụng của lực ma sát.
\choiceTF[t]
{\True Đế dép, lốp ô tô, lốp xe đạp phải khứa ở mặt cao su để tăng ma sát trượt}
{\True Việc sử dụng các con lăn, bánh xe, ổ bi trong kĩ thuật và đời sống là giảm tác hại của lực ma sát trượt}
{Cán cuốc khô khó cầm hơn cán cuốc ẩm ướt vì khi cán cuốc ẩm, các thớ gỗ phồng lên, ma sát tăng lên dễ cầm hơn}
{\True Khi dùng giấy nhám để đánh bóng bề mặt gỗ, người ta đã lợi dụng ma sát để làm nhẵn bề mặt}
\loigiai{
\begin{itemchoice}
\itemch
\itemch
\itemch
\itemch	
\end{itemchoice}
}
\haidong{10} \end{ex}
\begin{ex}%Tailieuchuan
Cho một vật có khối lượng $10\ kg$ đặt đứng yên trên sàn nhà. Tác dụng vào vật một lực có độ lớn $30\ N$ vật theo phương ngang. Hệ số ma sát giữa vật và sàn nhà là $\mu=0,2$. Cho $g=10 \ m/s^2$.
\choiceTF[t]
{\True Lực ma sát trượt tác dụng lên vật có độ lớn là $20\ N$} 
{Gia tốc của vật có độ lớn là $0,5 \ m/s^2$} 
{\True Tốc độ của vật sau $10\ s$ kể từ khi tác dụng lực là $10 \ m/s$} 
{Quãng đường vật đi được sau $10\ s$ kể từ khi tác dụng lực là $20\ m$} 
\loigiai{ 
\begin{itemchoice}
\itemch
\itemch
\itemch
\itemch
\end{itemchoice}
}
\end{ex}
\begin{ex}%Tailieuchuan
Một vật trượt từ đỉnh mặt phẳng nghiêng dài $10\ m$, cao $5\ m$. Bỏ qua ma sát trên mặt phẳng nghiêng. Sau khi đến chân mặt phẳng nghiêng, vật tiếp tục chuyển động trên mặt phẳng ngang trước khi dừng lại. Biết hệ số ma sát giữa vật và mặt phẳng ngang là $0,1$. Lấy $g=10 \ m/s^2$
\choiceTF[t]
{Gia tốc của vật trượt trên mặt phẳng nghiêng là $4 \ m/s^2$} 
{\True Tốc độ của vật khi trượt xuống đến chân dốc là $10 \ m/s$} 
{\True Gia tốc của vật khi trượt trên mặt phẳng ngang có độ lớn là $1 \ m/s^2$} 
{\True Thời gian vật chuyển động trên mặt phẳng ngang là $10\ s$} 
\loigiai{ 
\begin{itemchoice}
\itemch
\itemch
\itemch
\itemch
\end{itemchoice}
}
\end{ex}























\subsubsection*{PHẦN 3. Thí sinh trả lời từ câu 1 đến câu 6.}
\setcounter{ex}{0} 
\begin{ex}%book%SBT KNTT 18.4
Một vận động viên môn khúc côn cầu dùng gậy gạt bóng để truyền cho nó một tốc độ ban đầu $10\ m/s$. Hệ số ma sát giữa bóng và mặt băng là $0,1$. Lấy $g=9,8\ m/s^2$. Quãng đường quả bóng đi được cho đến khi dừng lại là bao nhiêu mét? 	
\shortans[oly]{51}
\haidong{10} 
\end{ex}
\begin{ex}%book%SBT KNTT 18.12
Một khúc gỗ khối lượng 2,5 kg đặt trên sàn nhà. Người ta kéo khúc gỗ bằng một lực $\vec F$ hướng chếch lên và hợp với phương nằm ngang một góc $\alpha=30^\circ$. Khúc gỗ chuyển động nhanh dần đều với gia tốc $1,5\ m/s^2$ trên sàn. Biết hệ số ma sát trượt giữa gỗ và sàn là 0,25. Lấy $g=10\ m/s^2$. Độ lớn của lực $\vec F$ là bao nhiêu niutơn (làm tròn kết quả đến chữ số hàng phần mười)?
\shortans[oly]{10,1}
\haidong{10}\loigiai{
\immini{\begin{itemize}
\item Ta có:
\begin{align*}
&F\cos\alpha-F_{mst}=ma\\
\Rightarrow &F\cos\alpha-\mu\left(P-F\sin\alpha\right)=ma\\
\Rightarrow &F=\dfrac{ma+\mu P}{\cos\alpha+\mu \sin\alpha}\approx 10,1\ N
\end{align*}
\end{itemize}}{\begin{tikzpicture}[scale=1,thick,>=stealth']
\tkzDefPoints{0/0.32/o}
\node  at  (o)  [shape=rectangle,draw=red,pattern=dots,minimum height=0.6cm,minimum width=1.4cm, rounded corners=1pt]  {};
\draw(-1.5,0)--(1.5,0);
\draw[very thick,blue,->](-0.7,0)--++(180:1.2)node[above,black]{$\vec{F}_{ms}$};
\draw[fill=white](-0.7,0)circle(1.2pt);
\draw[very thick,blue,->](o)--++(30:2)node[right,black]{$\vec{F}$}coordinate(f);
\draw[fill=white](o)circle(1.2pt);
\draw[very thick,blue,->](o)--++(0:1.73)node[below right,black]{$\vec{F}_1$}coordinate(ft);
\draw[very thick,blue,->](o)--++(90:1)node[left,black]{$\vec{F}_2$}coordinate(fn);
\draw[very thick,blue,->](o)--++(90:1.8)node[left,black]{$\vec{N}$};
\draw[very thick,blue,->](o)--++(-90:1.8)node[left,black]{$\vec{P}$};
\draw[dashed](ft)--(f)--(fn);
\fill[pattern=north east lines](-1.5,0)--(1.5,0)--++(-90:0.1)--++(180:3);
\draw[fill=white](o)circle(1.2pt);
\tkzLabelAngles[pos=1](ft,o,f){$\alpha$}
\end{tikzpicture}}
}
\haidong{10} \end{ex}





\begin{ex}%book%SBT KNTT 18.11
Một vật có khối lượng $2000\ g$ được đặt trên một bàn dài nằm ngang. Tác dụng lên vật một lực có độ lớn $5\ N$ theo phương song song với mặt bàn trong khoảng thời gian $2\ s$ rồi thôi tác dụng lực. Biết hệ số ma sát giữa vật và mặt bàn là $0,2$. Lấy $g=10\ m/s^2$. Quãng đường tổng cộng mà vật đi được cho đến khi dừng lại là bao nhiêu mét?
\shortans[oly]{1,25}
\haidong{10}\loigiai{
\begin{itemize}
\item Khi tác dụng lực: $F=ma+\mu mg\Rightarrow a=\dfrac{F-\mu mg}{m}=\dfrac{5-0,2.2.10}{2}=0,5\ m/s^2$.
\item Quãng đường vật đi được trong 2 s đầu:  $s_1=\dfrac{1}{2}at^2=\dfrac{1}{2}.0,5.2^2=1\ m$.
\item Sau 2 giây, vật chuyển động chậm dần dưới tác dụng của lực ma sát: $F_{ms}=ma'\Rightarrow a'=-\mu g=-2\ m/s^2$.
\item Quãng đường đi được từ lúc ngừng lực tác dụng tới khi dừng hẳn: $s_2=\dfrac{-v^2}{2a'}=\dfrac{-(at)^2}{2.(-2)}=0,25\ m$
\item Tổng quãng đường: $s=s_1+s_2=1,25\ m$
\end{itemize}
}
\haidong{10} \end{ex}
\begin{ex}%book
Cho một mặt phẳng nghiêng một góc $30^\circ$ so với phương ngang và có chiều dài $25\ m$. Đặt một vật tại đỉnh mặt phẳng nghiêng rồi cho trượt xuống thì có vận tốc ở cuối chân dốc là $10\ m/s$. Lấy $g=10\ m/s^2$. Hệ số ma sát giữa vật và mặt phẳng nghiêng là bao nhiêu (làm tròn kết quả đến chữ số hàng phần trăm)?
\shortans[oly]{0,35}
\haidong{10}
\loigiai{

\begin{itemize}
\item Có ba lực tác dụng lên vật khi vật trượt xuống mặt phẳng nghiêng: $\overrightarrow{N}, \overrightarrow{P}, \overrightarrow{F_{ms}}$

\item Áp dụng định luật II Niuton, ta có: $\overrightarrow{N}+\overrightarrow{P}+\overrightarrow{F_{ms}}=ma$

\item Chọn hệ trục gồm: Ox hướng theo chiều chuyển động của vật: trên mặt phẳng nghiêng, $Oy$ vuông góc với $Ox$ và hướng xuống.

\item Chiếu biểu thức vector (1) lên trục Ox, Oy ta được:

\end{itemize}

\begin{center}
\includegraphics[width=4cm]{images/2023_11_19_be4543c46a9f67e78874g-19}
\end{center}

Theo trục Ox: $P_{x}-F_{m s}=m a \Leftrightarrow P_{x}-\mu. N=m a(2)$

Theo trục Oy: $P_{y}-N=0$ (3)

Gia tốc của vật trên mặt phẳng nghiêng

$a=\dfrac{v^2-v_0^2}{2\ s}=\dfrac{10^2-0^2}{2.25}=2\ m/s^2$

Từ (1); (2); (3): ta có:

$\mu=\tan \alpha-\dfrac{a}{g \cos \alpha}=\dfrac{1}{\sqrt{3}}-\dfrac{2}{10 \cdot \dfrac{\sqrt{3}}{2}}=0,34$

Chọn A
}
\haidong{10} \end{ex}


\textit{Sử dụng thông tin sau cho Câu 5 và Câu 6:}
Một vật có khối lượng $m = 40$ kg bắt đầu trượt trên sàn nhà nằm ngang dưới tác dụng của một lực nằm ngang có độ lớn $200\ N$. Hệ số ma sát trượt giữa vật và sàn nhà là $0,25$. 
\begin{ex}%book%[1P2G7-2][Loai1]
Tốc độ của vật sau 5 giây kể từ khi bắt đầu trượt là bao nhiêu m/s?
\shortans[oly]{12,5}
\end{ex}

\begin{ex}%book%[1P2G7-2][Loai1]
Quãng đường đi được sau $5$ giây kể từ khi bắt đầu trượt là bao nhiêu mét (làm tròn kết quả đến chữ số hàng đơn vị)?
\shortans[oly]{31}
\loigiai{
\begin{itemize}
\item Các lực tác dụng lên vật là: $\vec F,\ {\vec F _{ms}}, N,\ \vec N,\ \vec P$
\item Chiếu lên phương chuyển động: $F - F_{ms} = ma \Rightarrow F - \mu mg =ma$ \\
$\Rightarrow a = \dfrac{{F - \mu mg}}{m} = \dfrac{{200 - 0,25.40.10}}{{40}} = 2,5\ m/s^2$
\item Vận tốc sau 5 giây là: $v = at = 2,5.5 = 12,5$ m/s
\item Quãng đường đi được sau 5 s là: $s = \dfrac{1}{2}at^2 = \dfrac{1}{2}.2,{5.5^2} = 31,25$ m.
\end{itemize}
}
\haidong{10} \end{ex}









\newpage 

\TieudeT{PHONG TỎA VẬT LÍ}
\subsubsection*{PHẦN I. Thí sinh trả lời câu hỏi từ câu 1 đến câu 18. Mỗi câu hỏi thí sinh chỉ chọn một phương án}
\setcounter{ex}{0}
\begin{ex}%book%[1P2B7-1][Loai1]
Đặc điểm nào sau đây phù hợp với lực ma sát trượt?
\choice
{\True Lực luôn xuất hiện ở mặt tiếp xúc và có hướng ngược với hướng chuyển động của vật}
{Lực luôn xuất hiện khi có sự biến dạng của vật}
{Lực xuất hiện khi có ngoại lực tác dụng vào vật nhưng vẫn đứng yên}
{Lực xuất hiện khi vật đặt gần bề mặt Trái Đất}
\loigiai{
Lực ma sát trượt luôn xuất hiện ở mặt tiếp xúc và có hướng ngược với hướng chuyển động của vật.
}
\haidong{10} \end{ex}
\begin{ex}%book%[1P2Y7-1][Loai1]
Lực ma sát trượt
\choice
{tỉ lệ thuận với vận tốc của vật}
{\True phụ thuộc vào độ lớn của áp lực}
{phụ thuộc vào diện tích mặt tiếp xúc}
{chỉ xuất hiện khi vật đang chuyển động chậm dần}
\loigiai{
Lực ma sát trượt xuất hiện khi vật này trượt lên vật khác, và có độ lớn bằng $F_{mst} = \mu N$ nên tỉ lệ thuận với áp lực.
}
\haidong{10} \end{ex}
\begin{ex}%book%SBT KNTT 18.1
Một vật đang trượt trên một mặt phẳng, khi tốc độ của vật giảm thì hệ số ma sát giữa vật và mặt phẳng
\choice
{giảm xuống}
{\True không đổi}
{tăng tỉ lệ với tốc độ của vật}
{tăng tỉ lệ với bình phương tốc độ của vật}
\haidong{10} \end{ex}
\begin{ex}%book%[1P2Y7-1][Loai1]
Nếu diện tích hai mặt tiếp xúc tăng lên thì hệ số ma sát trượt giữa hai mặt tiếp xúc
\choice
{tăng lên}
{giảm đi}
{tăng lên đến một giá trị cực đại}
{\True gần như không thay đổi}
\loigiai{
Hệ số ma sát hầu như không phụ thuộc vào diện tích mặt tiếp xúc mà phụ thuộc vào tính chất của các mặt tiếp xúc (có nhẵn hay không, làm bằng vật liệu gì).
}
\haidong{10} \end{ex}
\begin{ex}%book
Hệ số ma sát trượt
\choice
{không phụ thuộc vào vật liệu và tình trạng của hai mặt tiếp xúc}
{luôn nhỏ hơn hệ số ma sát lăn}
{\True không có đơn vị}
{tỉ lệ thuận với áp lực của vật lên mặt tiếp xúc}
\loigiai{
Hệ số ma sát trượt
\begin{itemize}
\item chỉ phụ thuộc vào tình trạng, chất liệu của hai mặt tiếp xúc;
\item không phụ áp lực của vật lên mặt tiếp xúc; tốc độ của vật; diện tích tiếp xúc giữa hai mặt tiếp xúc.
\item không có đơn vị
\item lớn hơn hệ số ma sát lăn nếu hai mặt tiếp xúc giống nhau.
\end{itemize}
}
\haidong{10} \end{ex}
\begin{ex}%book
Độ lớn của lực ma sát trượt tác dụng lên vật sẽ giảm $2$ lần nếu
\choice
{tăng hệ số ma sát lên $2$ lần}
{tăng diện tích tiếp xúc giữa vật và mặt tiếp xúc lên hai lần}
{giảm tốc độ chuyển động của vật $2$ lần}
{\True giảm áp lực vật đè lên mặt tiếp xúc đi $2$ lần}
\loigiai{
Lực ma sát trượt
\begin{itemize}
\item phụ thuộc vào tình trạng, chất liệu của hai mặt tiếp xúc;		
\item tỉ lệ với áp lực của vật lên mặt tiếp xúc;			
\item không phụ thuộc vào tốc độ của vật;			
\item không phụ thuộc vào diện tích tiếp xúc giữa hai mặt tiếp xúc.
\end{itemize}		
Nên nếu giảm áp lực 2 lần thì lực ma sát trượt giảm 2 lần
}
\haidong{10} 
\end{ex}
\begin{ex}%book%[1P2B7-1][Loai1]
Chiều của lực ma sát nghỉ
\choice
{ngược chiều với vận tốc của vật}
{ngược chiều với gia tốc của vật}
{\True ngược chiều với thành phần ngoại lực song song với mặt tiếp xúc}
{vuông góc với mặt tiếp xúc}
\loigiai{
Lực ma sát nghỉ ngược chiều với thành phần ngoại lực song song với mặt tiếp xúc.
}
\haidong{10} \end{ex}
\begin{ex}%book
Phát biểu nào sau đây \textbf{không} đúng khi nói về lực ma sát nghỉ?
\choice
{Lực ma sát nghỉ luôn xuất hiện ở bề mặt tiếp xúc giữa hai vật}
{Lực ma sát nghỉ giữ cho các điểm tiếp xúc của vật không trượt trên bề mặt}
{\True Một vật có thể tự đứng yên trên mặt phẳng nghiêng mà không cần đến lực ma sát nghỉ}
{Một vật có thể tự đứng yên trên mặt phẳng ngang mà không cần đến lực ma sát nghỉ}
\haidong{10} \end{ex}
\begin{ex}%book
Phát biểu nào sau đây là \textbf{sai}? Lực ma sát nghỉ
\choice
{có hướng ngược với hướng của thành phần lực song song với mặt tiếp xúc}
{có độ lớn bằng độ lớn của thành phần lực song song với mặt tiếp xúc}
{có phương song song với mặt tiếp xúc}
{\True là một lực luôn có hại}
\loigiai{
Lực ma sát nghỉ

\begin{itemize}
\item xuất hiện khi có một thành phần lực song song với mặt tiếp xúc tác dụng lên vật làm vật có xu hướng chuyển động nhưng vật chưa chuyển động

\item có phương song song với mặt tiếp xúc;

\item có hướng ngược với hướng của thành phần lực song song với mặt tiếp xúc.;

\item có độ lớn bằng độ lớn của thành phần lực song song với mặt tiếp xúc;

\item có nhiều lợi ích trong cuộc sống: giúp con người cầm nắm mọi vật, bức trang, đồ vật có thể treo trên tường, là lực phát động giúp con người chuyển động,....

\end{itemize}

Chọn D

}
\haidong{10} \end{ex}
\begin{ex}%book%SBT KNTT 18.2
Khi lực ép giữa hai mặt tiếp xúc tăng lên thì hệ số ma sát giữa hai mặt tiếp xúc	
\choice
{giảm đi}
{tăng lên}
{\True không thay đổi}
{không xác định được}
\loigiai{Hệ số ma sát chỉ phụ thuộc vào vật liệu và tình trạng bề mặt tiếp xúc}
\haidong{10} \end{ex}
\begin{ex}%book%[1P2Y7-1][Loai1]
Một vật lúc đầu nằm trên một máng nằm ngang. Sau khi được truyền một vận tốc đầu, vật chuyển động chậm dần vì có
\choice
{quán tính}
{phản lực}
{\True lực ma sát}
{lực tác dụng ban đầu}
\loigiai{
Sau khi được truyền một vận tốc đầu, vật chuyển động chậm dần vì có lực ma sát.
}
\haidong{10} 
\end{ex}
\begin{ex}%book%[1P2B7-1][Loai1]
Thủ môn bắt "dính" bóng là nhờ
\choice
{lực ma sát trượt}
{\True lực ma sát nghỉ}
{lực quán tính}
{lực ma sát lăn}
\loigiai{
Thủ môn bắt "dính" bóng là nhờ lực ma sát nghỉ.
}
\haidong{10} \end{ex}
\begin{ex}%book%SBT KNTT 18.5
Một xe tải có khối lượng $3$ tấn đang chuyển động trên đường nằm ngang, hệ số ma sát của xe tải với mặt đường là $0,1$. Lấy $g=10\ m/s^2$. Độ lớn của lực ma sát là
\choice
{\True $3000\ N$}
{$30000\ N$}
{$300\ N$}
{$30\ N$}
\haidong{3}
\haidong{10} \end{ex}
\begin{ex}%book
Khi đẩy một ván trượt bằng một lực $F_1=20\ N$ theo phương ngang thì nó chuyển động thẳng đều. Nếu chất lên ván một hòn đá nặng $20\ kg$ thì để nó trượt đều phải tác dụng lực $F_2=$ $60\ N$ theo phương ngang. Hệ số ma sát trượt giữa tấm ván và mặt sàn là
\choice
{$0,25$}
{\True $0,2$}
{$0,1$}
{$0,15$}
\loigiai{
Đặt $m$ là khối lượng tấm ván, $\Delta m$ là khối lượng hòn đá. Do cả hai trường hợp đều trượt đều $(a=0)$ nên ta có:

$$
\begin{aligned}
& \left\{\begin{array}{l}
F_1=\mu m g \\
F_2=\mu(m+\Delta m) g
\end{array} \Rightarrow F_2-F_1=\mu \cdot \Delta m \cdot g\right. \\
& \Rightarrow \mu=\dfrac{F_2-F_1}{\Delta m \cdot g}=0,2
\end{aligned}
$$
}
\haidong{10} \end{ex}
\begin{ex}%book%[1P2K7-2][Loai1]
Vật có khối lượng $2\ kg$ đặt trên mặt bàn nằm ngang. Hệ số ma sát trượt giữa vật và bàn là $\mu = 0,25$. Tác dụng một lực $6\ N$ lên vật theo phương song song với mặt bàn. Cho $g = 10\ m/s^2$. Gia tốc của vật chuyển động có độ lớn là
\choice
{$0,2\ m/s^2$}
{\True $0,5\ m/s^2$}
{$1,0\ m/s^2$}
{$1,2\ m/s^2$}
\haidong{4}\loigiai{
Gia tốc chuyển động của vật là: $a = \dfrac{F - F_{ms}}{m} = \dfrac{F - \mu mg}{m} =\dfrac{6 - 0,25.2.10}{2} = 0,5\ m/s^2$.
}
\haidong{10} \end{ex}
%ldd
\begin{ex}%book%[1P2K7-2][Loai1]
Người ta đẩy một cái thùng có khối lượng $55\ kg$ theo phương ngang với lực $220\ N$ làm thùng chuyển động trên mặt phẳng ngang. Hệ số ma sát trượt giữa thùng và mặt phẳng là $0,35$. Lấy $g = 9,8\ m/s^2$. Gia tốc của thùng có độ lớn là
\choice
{$0,43\ m/s^2$}
{$0,35\ m/s^2$}
{\True $0,57\ m/s^2$}
{$0,6\ m/s^2$}
\haidong{4}\loigiai{
Gia tốc chuyển động của thùng là: $a = \dfrac{F - F_{ms}}{m} = \dfrac{220 - 0,35.55.9,8}{55} = 0,57\ m/s^2$.
}
\haidong{10} \end{ex}
\begin{ex}%book%SBT KNTT 18.10
Một vật có khối lượng $15\ kg$ đang đứng yên trên mặt phẳng ngang thì bắt đầu chuyển động nhanh dần đều dưới tác dụng của lực kéo theo phương ngang, sau khi đi được $150\ m$ vật đạt tốc độ $54\ km/h$. Biết hệ số ma sát giữa vật và mặt phẳng ngang là $0,05$. Lấy $g=9,8\ m/s^2$. Lực kéo tác dụng vào vật có độ lớn là
\choice
{$16,8\ N$}
{$28,6\ N$}
{$23,7\ N$}
{\True $18,6\ N$}
\haidong{5}\loigiai{
\begin{itemize}
\item $F=ma+\mu mg=m\dfrac{v^2}{2s}+\mu mg=18,6\ N$.
\end{itemize}
}
\haidong{10} 
\end{ex}
\begin{ex}%book%[1P2G7-2][Loai1]
\immini[thm]{Một cái hòm có khối lượng $m = 20\ kg$ đặt trên sàn nhà. Người ta kéo hòm bằng một lực $F$ hướng chếch lên trên và hợp với phương nằm ngang một góc $\alpha = 20^\circ$ như hình vẽ. Hòm chuyển động đều trên sàn nhà. Hệ số ma sát trượt giữa hòm và sàn nhà $\mu_t = 0,3$. Độ lớn của lực $F$ là }{\begin{tikzpicture}[scale=1,thick,>=stealth']
\tkzDefPoints{0/1/o}
\node at (0,0.32) [shape=rectangle,draw=red,pattern=dots,minimum height=0.6cm,minimum width=1.4cm, rounded corners=1pt] {};
\draw(-1.5,0)--(1.5,0);
\draw[very thick,blue,->](0.7,0.6)--++(30:2)node[right,black]{$\vec{F}$}coordinate(f);
\draw[fill=white](0.7,0.6)circle(1.2pt);
\fill[pattern=north east lines](-1.5,0)--(1.5,0)--++(-90:0.1)--++(180:3);
\end{tikzpicture}}
\choice
{$52,3\ N$}
{$46,5\ N$}
{$42,6\ N$}
{\True $56,4\ N$}
\haidong{10}\loigiai{
\immini{Để vật chuyển động thẳng đều thì $F_{ms} = F\cos \alpha \Rightarrow \left( {P - F\sin \alpha } \right) = P\cos \alpha$\\ $\Rightarrow F = \dfrac{{\mu P}}{{\mu \sin \alpha + \cos \alpha }} = \dfrac{{0,3.20.10}}{{0,3.\sin {{20}^\circ} + \cos {{20}^\circ}}} = 56,4\ N.$}{\begin{tikzpicture}[scale=1,thick,>=stealth']
\tkzDefPoints{0/1/o}
\node at (0,0.32) [shape=rectangle,draw=red,pattern=dots,minimum height=0.6cm,minimum width=1.4cm, rounded corners=1pt] {};
\draw(-1.5,0)--(1.5,0);
\draw[very thick,blue,->](-0.7,0)--++(180:1.2)node[above,black]{$\vec{F}_{ms}$};
\draw[fill=white](-0.7,0)circle(1.2pt);
\draw[very thick,blue,->](0.7,0.6)--++(30:2)node[right,black]{$\vec{F}$}coordinate(f);
\draw[fill=white](0.7,0.6)circle(1.2pt);
\draw[very thick,blue,->](0.7,0.6)--++(0:1.73)coordinate(ft);
\draw[very thick,blue,->](0.7,0.6)--++(90:1)coordinate(fn);
\draw[dashed](ft)--(f)--(fn);
\fill[pattern=north east lines](-1.5,0)--(1.5,0)--++(-90:0.1)--++(180:3);
\draw[very thick,blue,->](0,0.32)--++(90:1.8)node[left,black]{$\vec{R}$};
\draw[very thick,blue,->](0,0.32)--++(-90:1.8)node[left,black]{$\vec{P}$};
\draw[fill=white](0,0.32)circle(1.2pt);
\end{tikzpicture}}	
}
\haidong{10} \end{ex}




\subsubsection*{PHẦN 2. Thí sinh trả lời từ câu 1 đến câu 4. Trong mỗi ý a), b), c), d) ở mỗi câu, thí sinh chọn đúng hoặc sai.}
\setcounter{ex}{0}
\begin{ex}%Tailieuchuan
Một xe đạp đang chuyển động trên đường.
\choiceTF[t]
{Nếu xe đang chuyển động thẳng đều  thì lực ma sát giữa bánh xe và mặt đường không tồn tại} 
{\True Lực ma sát giữa bánh xe và mặt đường có tác dụng giúp xe chuyển động} 
{\True Khi xe phanh gấp thì lực ma sát trượt giữa bánh xe và mặt đường giúp xe dừng lại} 
{Diện tích tiếp xúc giữa bánh xe và mặt đường càng lớn thì lực ma sát càng giảm} 
\loigiai{ 
\begin{itemchoice}
\itemch
\itemch
\itemch
\itemch
\end{itemchoice}
}
\end{ex}
\begin{ex}%Tailieuchuan
Một cuốn sách trượt trên mặt bàn phẳng bàn nằm ngang khi được truyền vận tốc đầu.
\choiceTF[t]
{Lực ma sát trượt giữa cuốn sách và mặt bàn giúp cuốn sách tăng tốc} 
{Khi cuốn sách trượt trên mặt phẳng, lực ma sát trượt giữa cuốn sách và mặt bàn phụ thuộc vào tốc độ trượt của cuốn sách} 
{\True Cuốn sách sẽ dừng lại sau một thời gian do lực ma sát trượt giữa sách và mặt bàn} 
{\True Độ lớn của lực ma sát trượt phụ thuộc vào trọng lượng của cuốn sách} 
\loigiai{ 
\begin{itemchoice}
\itemch
\itemch
\itemch
\itemch
\end{itemchoice}
}
\end{ex}
\begin{ex}%Tailieuchuan
Một người kéo một chiếc vali trên nền nhà bằng tay cầm.
\choiceTF[t]
{Khi vali đang di chuyển đều, lực ma sát giữa vali và mặt nền bằng không} 
{\True Lực kéo tác dụng lên vali tối thiểu phải lớn hơn lực ma sát thì vali mới có thể chuyển động} 
{\True Lực ma sát giữa bánh xe của vali và mặt nền phụ thuộc vào độ nhám của bánh xe và mặt nền} 
{Diện tích tiếp xúc giữa bánh xe và mặt nền càng lớn thì lực ma sát càng lớn} 
\loigiai{ 
\begin{itemchoice}
\itemch
\itemch
\itemch
\itemch
\end{itemchoice}
}
\end{ex}
\begin{ex}%Tailieuchuan
Một ô tô có khối lượng $3,6$ tấn bắt đầu chuyển động trên đường nằm ngang với lực kéo $F$ theo phương ngang. Sau $20\ s$ vận tốc của xe là $15 \ m/s$. Biết lực ma sát của xe với mặt đường bằng $0,25F$. Lấy $g=10\ m/s^2$.
\choiceTF[t]
{\True Gia tốc của vật có độ lớn là $0,75\ m/s^2$}
{Lực kéo tác dụng lên vật có độ lớn là $F=360\ N$} 
{Lực ma sát tác dụng lên vật có độ lớn là $F_{ms}=90\ N$} 
{\True Hệ số ma sát giữa vật và mặt đường là $\mu=0,0252$} 
\loigiai{ 
\begin{itemchoice}
\itemch
\itemch
\itemch
\itemch
\end{itemchoice}
}
\end{ex}



















\subsubsection*{PHẦN 3. Thí sinh trả lời từ câu 1 đến câu 6.}
\setcounter{ex}{0} 
\begin{ex}%Tailieuchuan
Một vật khối lượng $400\ g$ đang trượt đều trên mặt phẳng ngang dưới tác dụng của lực $F=0,4\ N$. Lấy $g=10\ m/s^2$. Hệ số ma sát trượt giữa vật và mặt phẳng ngang là bao nhiêu?
\shortans[oly]{0,1}
\loigiai{}
\end{ex}
\begin{ex}%Tailieuchuan
\immini[thm]{Cho một vật có khối lượng $100\ kg$ đang đứng yên trên mặt phẳng nằm ngang. Tác dụng vào vật một lực $F$  $\alpha=$ $30^{0}$, Hệ số ma sát $\mu=0,2$ và $g=10\ m/s^2$. Để vật chuyển động đều thì độ lớn của lực là bao nhiêu niutơn (làm tròn kết quả đến chữ số hàng đơn vị)?}{\begin{tikzpicture}[draw=Mapcolor, very thick,scale=1,>=stealth',draw=Mapcolor]
\fill[pattern=north east lines] (0,-0.2) rectangle (6,0);
\draw (0,0)--(6,0);
\fill[pattern=north west lines,draw] (2,0) rectangle (3.5,0.8);
\draw[->,very thick,Mapcolor](3.5,0.4)--++(30:2)node[right]{$\vec F$};
\draw[densely dashed] (3.5,0.4)--(6,0.4);
\draw (3.5,0.4)++(0:0.6) arc (0:30:0.6);
\path(3.5,0.4)--++(15:0.8) node{{$\alpha$}};
\end{tikzpicture}}
\shortans[oly]{207}
\loigiai{
}
\end{ex}
\textit{Sử dụng thông tin sau cho Câu 3 và Câu 4:} Một mẩu gỗ có khối lượng $m = 250\ g$ đặt trên sàn nhà nằm ngang. Người ta truyền cho nó một vận tốc tức thời $v_0 = 5\ m/s$. Hệ số ma sát trượt giữa mẩu gỗ và sàn nhà là $\mu_t = 0,25$. Lấy $g = 10\ m/s^2$.  \begin{ex}%book%[1P2G7-2][Loai1]
Thời gian để mẩu gỗ dừng lại là bao nhiêu giây?
\shortans[oly]{2}
\haidong{10} \end{ex}
\begin{ex}%book%[1P2G7-2][Loai1]
Quãng đường nó đi được cho tới khi dừng lại là bao nhiêu mét?
\shortans[oly]{5}	
\haidong{10}
\loigiai{
\begin{itemize}
\item Ta có: $F_{ms} = \mu P = \mu mg \Rightarrow a = \dfrac{- F_{ms}}{m}=-\mu g=- 2,5\ m/s^2$
\item Áp dụng công thức độc lập thời gian có: $v^2 - v_0^2 = 2as$ 
\item Quãng đường bóng đi được là: $s = \dfrac{v^2 - v_0^2}{2a} = \dfrac{0^2 - 5^2}{2.\left(-2,5\right)} = 5\ m$ 
\item Thời gian mẩu gỗ chuyển động là: $t = \dfrac{v - v_0}{a} = \dfrac{0 - 5}{-2,5} = 2\ s$.
\end{itemize}
}
\haidong{10} \end{ex}

\begin{ex}%book%[1P2K7-2][Loai1]
Hùng và Dũng cùng nhau đẩy một thùng hàng chuyển động thẳng trên sàn nhà nằm ngang. Thùng hàng có khối lượng $120\ kg$. Hùng đẩy với một lực $400\ N$, Dũng đẩy với một lực $300\ N$, các lực đều có phương song song với sàn nhà. Hệ số ma sát trượt giữa thùng và sàn nhà là $\mu_t = 0,2$. Lấy $g = 10\ m/s^2$. Gia tốc của thùng hàng có độ lớn là bao nhiêu $m/s^2$ (làm tròn kết quả đến chữ số hàng phần mười)?
\shortans[oly]{3,8}
\haidong{10}\loigiai{
\immini{\begin{itemize}
\item Các lực tác dụng lên thùng hàng như hình vẽ. 
\item Trong đó $\vec F$ là tổng hợp lực đẩy của Hùng và Dũng tác dụng vào thùng hàng, có độ lớn là $F = 400 + 300 = 700$ N.
\item Ta có: $F_{ms} = \mu mg = 0,2.120.10 = 240$ N
\item Gia tốc chuyển động của thùng là $a = \dfrac{F - F_{ms}}{m} = \dfrac{{700 - 240}}{120} = 3,8\ m/s^2$.
\end{itemize}}{\begin{tikzpicture}[scale=1,thick,>=stealth']
\tkzDefPoints{0/1/o}
\node at (0,0.32) [shape=rectangle,draw=red,pattern=dots,minimum height=0.6cm,minimum width=1.4cm, rounded corners=1pt] {};
\draw(-1.5,0)--(1.5,0);
\draw[very thick,blue,->](0.1,0)--++(90:1.8)node[right,black]{$\vec{N}$};
\draw[fill=white](0.1,0)circle(1.2pt);
\draw[very thick,blue,->](0,0.32)--++(-90:1.8)node[right,black]{$\vec{P}$};
\draw[fill=white](0,0.32)circle(1.2pt);
\draw[very thick,blue,->](-0.7,0)--++(180:1.2)node[above,black]{$\vec{F}_{ms}$};
\draw[fill=white](-0.7,0)circle(1.2pt);
\draw[very thick,blue,->](0.7,0.32)--++(0:1.5)node[above,black]{$\vec{F}$};
\draw[fill=white](0.7,0.32)circle(1.2pt);
\fill[pattern=north east lines](-1.5,0)--(1.5,0)--++(-90:0.1)--++(180:3);
\end{tikzpicture}}
}
\haidong{10} \end{ex}

\begin{ex}%book%[1P2G7-2][Loai1]
\immini[thm]{Một vật nhỏ đặt trên một máng nghiêng $MN$ khá dài hợp với mặt phẳng nằm ngang một góc $\alpha = 20^\circ$. Hệ số ma sát nghỉ và hệ số ma sát trượt giữa vật và máng nghiêng đều có trị số là $\mu = 0,2$. Truyền cho vật một vận tốc ban đầu $v_0$ như hình vẽ. Lấy $g = 10\ m/s^2$. Gia tốc của vật có độ lớn là bao nhiêu $m/s^2$ (làm tròn kết quả đến chữ số hàng phần mười)?
\shortans[oly]{5,3}}{\begin{tikzpicture}[scale=1,thick,>=stealth']
\tkzInit[ymin=-4,ymax=0.5,xmin=-0.5,xmax=5.5]\tkzClip
\begin{scope}[shift={(0,0)},rotate=-30]
\node at (2.3,0) [shape=rectangle,draw=red,pattern=dots,minimum height=0.5cm,minimum width=0.6cm, rounded corners=1pt,rotate=-30] {};
\draw[very thick,blue,->](2,0)--++(180:0.5)node[above,black,xshift=-10pt]{$\vec{v}_{0}$};
\draw[fill=white](2,0)circle(1.2pt);	
\draw(0,-0.27)coordinate(a)--++(0:6)coordinate(b)--++(210:5)coordinate(c);
\fill[pattern=north east lines](b)--(c)--++(300:0.1)--++(390:5);
\tkzMarkAngles[size=0.5cm](a,b,c)
\tkzLabelAngle[pos=-0.8](a,b,c){$\alpha$}
\end{scope}
\end{tikzpicture}}
\haidong{10}\loigiai{
\immini{	Ta có: $F_{ms} = \mu N = \mu mg\cos \alpha$ $\Rightarrow$ $a = - \dfrac{\mu mg\cos \alpha + mg\sin \alpha}{m} = - g\left( {\mu \cos \alpha + \sin \alpha } \right) = -5,3\ m/s^2$}{\begin{tikzpicture}[scale=1,thick,>=stealth']
\tkzInit[ymin=-4,ymax=0.5,xmin=-0.5,xmax=5.5]\tkzClip
\begin{scope}[shift={(0,0)},rotate=-30]
\node at (2.3,0) [shape=rectangle,draw=red,pattern=dots,minimum height=0.5cm,minimum width=0.6cm, rounded corners=1pt,rotate=-30] {};
\draw[very thick,blue,->](2,0)--++(180:0.5)node[above,black,xshift=-10pt]{$\vec{v}_{0}$};
\draw[fill=white](2,0)circle(1.2pt);
\draw[very thick,blue,->](2.6,-0.25)--++(0:1)node[right,black]{$\vec{F}_{ms}$};
\draw[fill=white](2.6,-0.25)circle(1.2pt);
\draw[very thick,blue,->](2.3,0)--++(0:0.7)coordinate(px);
\draw[very thick,blue,->](2.3,0)coordinate(o')--++(90:1.2)node[right,black]{$\vec{R}$};
\draw[very thick,blue,->](2.3,0)--++(-90:1.2)coordinate(n);
\draw[very thick,blue,->](2.3,0)--++(-60:1.4)node[right,black,yshift=5pt]{$\vec{P}$}coordinate(p);
\draw[dashed,thin](px)--(p)--(n);
\draw[fill=white](2.3,0)circle(1.2pt);
\draw(0,-0.27)coordinate(a)--++(0:6)coordinate(b)--++(210:5)coordinate(c);
\fill[pattern=north east lines](b)--(c)--++(300:0.1)--++(390:5);
\tkzMarkAngles[size=0.5cm](a,b,c)
\tkzLabelAngle[pos=-0.8](a,b,c){$\alpha$}
\tkzMarkAngles[size=0.5cm](px,p,o')
\tkzLabelAngle[pos=0.7](px,p,o'){$\alpha$}
\end{scope}
\end{tikzpicture}}
}
\haidong{10} \end{ex}






